\section{Related Work}

\label{sec:related}
% ══════════════════════════════════════════════════════════════════════════════

\subsection{DAO Governance}

DAOs coordinate collective decision-making through smart contracts.
MakerDAO pioneered DeFi governance \citep{mcdao2017}; Compound's
Governor framework standardized on-chain proposal-vote-execute patterns
\citep{compound2020governance}. However, governance participation remains
low and plutocratic \citep{barbereau2022defi}.

\subsection{Quadratic Mechanisms}

Quadratic voting (QV) \citep{lalley2018quadratic} makes the cost of
expressing preference intensity scale quadratically: $K$ votes cost
$K^2$ credits. This produces welfare-optimal outcomes under standard
assumptions \citep{posner2018radical}. Quadratic funding (QF)
\citep{buterin2019flexible} extends QV to public goods, where matching
funds amplify small contributions. Gitcoin has deployed QF for Ethereum
ecosystem grants, distributing over \$50M \citep{gitcoin2023}.

\subsection{Reputation Systems}

On-chain reputation is implemented through soulbound tokens (SBTs)
\citep{weyl2022decentralized}---non-transferable tokens representing
credentials. SourceCred \citep{sourcecred2020} and Coordinape
\citep{coordinape2021} provide peer attestation-based alternatives.

\subsection{Science DAOs}

Existing science DAOs include VitaDAO (longevity), LabDAO (computational
biology), Molecule (drug discovery IP-NFTs), and ResearchHub (open science
discussion). Zoo DAO differs by focusing on conservation science,
incorporating multi-stakeholder governance, and implementing quadratic
mechanisms designed for research allocation.


% ══════════════════════════════════════════════════════════════════════════════
